\documentclass[../main.tex]{subfiles}
\usepackage[utf8]{inputenc}
\usepackage[T1]{fontenc}
\usepackage[indonesian]{babel}

\begin{document}

\chapter{Kriptografi Kurva Eliptik (ECC)}

\begin{learningoutcome}
Setelah mempelajari bab ini, mahasiswa diharapkan mampu:
\begin{itemize}
    \item Menjelaskan konsep dasar aljabar abstrak (grup, medan, dan Medan Galois) yang mendasari ECC (Sub-CPMK 2.3).
    \item Menganalisis properti geometris dan analitik dari kurva eliptik, termasuk operasi penjumlahan dan penggandaan titik (Sub-CPMK 2.3).
    \item Menguraikan mekanisme kerja \textit{Elliptic Curve Discrete Logarithm Problem} (ECDLP) sebagai basis keamanan ECC (Sub-CPMK 2.3).
    \item Menerapkan prosedur pertukaran kunci ECDH dan enkripsi ECEG (Sub-CPMK 2.3).
    \item Mengevaluasi efisiensi ECC dibandingkan dengan RSA dari sisi panjang kunci dan kebutuhan sumber daya komputasi (Sub-CPMK 2.3).
    \item Menjelaskan metode \textit{encoding} pesan menjadi titik pada kurva eliptik (Sub-CPMK 2.3).
\end{itemize}
\end{learningoutcome}

\section{Landasan Aljabar Abstrak}
Kriptografi Kurva Eliptik (ECC) dibangun di atas konsep aljabar abstrak untuk memastikan operasi matematika yang konsisten dan aman.

\subsection{Grup dan Medan}
\begin{itemize}
    \item \textbf{Grup}: Sebuah himpunan $G$ dengan operasi biner $*$ yang memenuhi aksioma \textit{closure}, asosiatif, memiliki elemen netral, dan setiap elemen memiliki invers. Jika berlaku komutatif ($a*b = b*a$), maka disebut \textbf{Grup Abelian}.
    \item \textbf{Medan (\textit{Field})}: Struktur aljabar $\langle F, +, \cdot \rangle$ di mana $\langle F, + \rangle$ adalah grup abelian, $\langle F - \{0\}, \cdot \rangle$ adalah grup abelian, dan operasi perkalian bersifat distributif terhadap penjumlahan.
\end{itemize}

\subsection{Medan Galois $\text{GF}(p^m)$}
ECC umumnya beroperasi pada medan berhingga atau \textbf{Medan Galois}, dinotasikan sebagai $\text{GF}(p^m)$.
\begin{itemize}
    \item \textbf{GF($p$)}: Medan dengan elemen $\{0, 1, \dots, p-1\}$ di mana operasi dilakukan modulo $p$.
    \item \textbf{GF($2^m$)}: Medan biner di mana elemen direpresentasikan sebagai polinomial dengan koefisien dalam $\text{GF}(2)$. Operasi perkalian dilakukan modulo \textit{irreducible polynomial}.
\end{itemize}

\section{Geometri Kurva Eliptik}
Kurva eliptik yang digunakan dalam kriptografi umumnya memiliki bentuk persamaan Weierstrass:
\[ y^2 = x^3 + ax + b \]
dengan syarat $4a^3 + 27b^2 \neq 0$ untuk memastikan kurva tidak memiliki titik singular (titik tajam atau perpotongan diri).

\subsection{Operasi Penjumlahan Titik}
Titik-titik pada kurva eliptik bersama dengan titik tak hingga $\mathcal{O}$ (elemen netral) membentuk sebuah grup abelian.
\begin{enumerate}
    \item \textbf{Penjumlahan ($P + Q = R$)}: Jika $P \neq Q$, tarik garis melalui $P$ dan $Q$. Garis tersebut akan memotong kurva di titik $-R$. Hasil penjumlahan $R$ adalah pencerminan $-R$ terhadap sumbu-x.
    \item \textbf{Penggandaan ($2P = R$)}: Jika $P = Q$, tarik garis tangen pada titik $P$. Titik potong garis tangen dengan kurva adalah $-R$, dan hasil akhirnya adalah $R$.
\end{enumerate}

\subsection{Analisis Koordinat ( Medan $\mathbb{R}$)
Untuk dua titik $P(x_p, y_p)$ dan $Q(x_q, y_q)$, koordinat hasil $R(x_r, y_r)$ dihitung dengan:
\[ m = \frac{y_q - y_p}{x_q - x_p} \quad \text{(untuk } P \neq Q \text{)} \quad \text{atau} \quad m = \frac{3x_p^2 + a}{2y_p} \quad \text{(untuk } P = Q \text{)} \]
\[ x_r = m^2 - x_p - x_q \quad \text{dan} \quad y_r = m(x_p - x_r) - y_p \]

\section{Kriptografi Berbasis Kurva Eliptik}
Inti dari keamanan ECC adalah \textbf{Elliptic Curve Discrete Logarithm Problem (ECDLP)}.

\subsection{ECDLP}
Diberikan titik basis $P$ dan titik $Q$ pada kurva eliptik, sangat sulit untuk menemukan integer $k$ sedemikian sehingga:
\[ Q = kP = \underbrace{P + P + \dots + P}_{k \text{ kali}} \]
Jika $k$ adalah bilangan yang sangat besar, mencari $k$ dari $P$ dan $Q$ secara komputasi tidak praktis.

\subsection{Mekanisme Kunci}
\begin{itemize}
    \item \textbf{Kunci Privat}: Integer acak $k \in [1, p-1]$.
    \item \textbf{Kunci Publik}: Titik $Q = kP$ pada kurva eliptik.
\end{itemize}

\subsection{Protokol Utama}
\begin{enumerate}
    \item \textbf{ECDH (Elliptic Curve Diffie-Hellman)}: Pertukaran kunci rahasia bersama menggunakan operasi perkalian titik.
    \item \textbf{ECEG (Elliptic Curve ElGamal)}: Enkripsi pesan dengan mengkodekan pesan menjadi titik $P_M$ dan menghasilkan cipherteks berupa pasangan titik $[(kB), (P_M + kP_B)]$.
\end{enumerate}

\section{Implementasi dan Efisiensi}
\subsection{Encoding Pesan}
Karena ECC beroperasi pada titik-titik kurva, pesan teks harus diubah menjadi koordinat $(x, y)$. Salah satu metodenya adalah menggunakan parameter basis $k$:
\[ x = mk + \text{offset} \]
Kemudian mencari nilai $y$ yang memenuhi persamaan kurva. Jika tidak ditemukan, offset ditambah hingga ditemukan titik yang valid.

\subsection{Perbandingan Keamanan: ECC vs RSA}
Kelebihan utama ECC adalah efisiensi panjang kunci. Untuk tingkat keamanan yang sama, ECC membutuhkan kunci yang jauh lebih pendek daripada RSA:
\begin{itemize}
    \item \textbf{Kunci AES-128} $\approx$ \textbf{Kunci RSA 3072-bit} $\approx$ \textbf{Kunci ECC 256-bit}.
\end{itemize}
Hal ini membuat ECC sangat ideal untuk perangkat dengan sumber daya terbatas seperti \textit{smart cards} dan perangkat IoT.

\begin{obeactivity}
\textbf{Aktivitas 14.1: Simulasi Penjumlahan Titik}
1. Gunakan kurva $y^2 = x^3 + 2x + 4$ di $\mathbb{R}$.
2. Diberikan titik $P(2, 4)$ dan $Q(0, 2)$, hitunglah koordinat titik $R = P + Q$.
3. Lakukan penggandaan titik untuk menghitung $2P$ pada kurva yang sama.
4. Bandingkan hasil perhitungan manual Anda dengan kalkulator ECC online.
\end{obeactivity}

\begin{obereflection}
\textbf{Latihan 14.1}
1. Mengapa titik tak hingga $\mathcal{O}$ diperlukan dalam definisi grup kurva eliptik?
2. Jelaskan mengapa ECDLP dianggap lebih sulit dipecahkan daripada masalah pemfaktoran bilangan bulat pada RSA.
3. Apa dampak penggunaan Medan Galois $\text{GF}(p)$ dibandingkan bilangan riil dalam implementasi praktis ECC?
\end{obereflection}

\begin{obeassessment}
\textbf{Kuis Bab 14}
1. Uraikan aksioma grup yang harus dipenuhi oleh titik-titik pada kurva eliptik.
2. Bagaimana cara menghitung gradien $m$ untuk operasi penggandaan titik $2P$?
3. Mengapa ECC lebih cocok untuk perangkat mobile dibandingkan RSA?
\end{obeassessment}

\begin{competencychecklist}
\begin{tabular}{|l|c|}
\hline
Indikator Kompetensi & Tercapai \\
\hline
Saya dapat menjelaskan konsep grup dan medan dalam ECC & [ ] \\
\hline
Saya dapat menghitung penjumlahan dan penggandaan titik pada kurva & [ ] \\
\hline
Saya dapat menguraikan mekanisme ECDLP & [ ] \\
\hline
Saya dapat membedakan ECDH dan ECEG & [ ] \\
\hline
Saya dapat menganalisis trade-off panjang kunci RSA vs ECC & [ ] \\
\hline
\end{tabular}
\end{center}
\end{competencychecklist}

\end{document}
